% 光的干涉（高中）
% keys 波的叠加
% license Usr
% type Tutor
\pentry{机械波及叠加（高中）\nref{nod_mcwave}}{nod_77a3}
在\autoref{sub_mcwave_1}，我们导出了如下规律，倘若两个简谐波的初相位差为0，那么空间中波程差为$n\lambda$的位置是振动加强点，波程差为$(2n+1)\lambda/2$则为振动减弱点。接下来，我们将利用这个规律求得干涉实验的一些特点。

首先，我们回顾一下光的构成。作为电磁波，有电场和磁场两个分量，且都在做简谐振动，振动方向与传播方向垂直，如下图所示。
\begin{figure}[ht]
\centering
\includegraphics[width=12cm]{./figures/add5a31d3e87b27a.png}
\caption{} \label{fig_interf_2}
\end{figure}
无论是人眼的光电化学反应，或是其他光电检测设备，所能检测的都是光的强度（即光的能流密度）。可以证明，光强正比于电场振幅在一段时间的平均值。又因为电场在做简谐振动，于是在

\subsection{双缝干涉}
\subsubsection{理论推导}

\subsubsection{实验装置}
一个经典的双缝干涉实验装置如下图所示。光先通过单缝，成为柱面波，我们可以画出一系列同心圆，每一个圆代表一个相位，最前面的波面是光源的初始相位。也就是说，柱面波到达双缝的必是光源的初始相位。由此，我们在双缝处得到两个相位差为0的光源。
\begin{figure}[ht]
\centering
\includegraphics[width=12cm]{./figures/17f92705aa639da5.png}
\caption{} \label{fig_interf_1}
\end{figure}

\subsection{薄膜干涉}